% --- Archon physics formalization source begin ---
% archon:physics
% archon:covers PhyXMiniProblems/problem_phyx_mini_0460.lean
% archon:source-report reports/phyx_mini/problem_phyx_mini_0460.source.json

\chapter{Physics problem phyx\_mini\_0460}
\label{ch:PhyXMiniProblems_problem_phyx_mini_0460}

\paragraph{Problem source.}
\#\# Physical scenario

A piston/cylinder assembly in a car contains \$0.2 \textbackslash{}, \textbackslash{}text\{L\}\$ of air at \$90 \textbackslash{}, \textbackslash{}text\{kPa\}\$ and \$20 \textbackslash{}, \textasciicircum{}\{\textbackslash{}circ\}\textbackslash{}text\{C\}\$, as shown in figure. The air is compressed in a quasi-equilibrium polytropic process with polytropic exponent \$n = 1.25\$ to a final volume six times smaller.

\#\# Question

Determine the heat transfer for the process

\#\# Answer choices

- A: A: -0.190 kJ

- B: B: 0.0147 kJ

- C: C: 0.190 kJ

- D: D: -0.0147 kJ

\#\# Figure caption (auxiliary; use the image as primary evidence)

The image is a schematic representation that appears to show a cross-sectional view of a closed system or apparatus with two main sections. On the left side, there is a light blue area labeled "Air," indicating it is filled with air. A solid vertical line separates this region from the right section.

In the right section, there are two large, looped shapes situated one above the other. The left loop is drawn with dashed lines, suggesting a boundary or a possible area of interest which may change or interact with the air. The right loop is a solid, continuous line and is positioned within a rectangular boundary that is connected to the rest of the structure.

The entire figure is enclosed within a larger rectangle which seems to represent the external structure or casing. The arrangement implies some kind of process or interaction between the outlined components, possibly involving airflow or ventilation.

\#\# Dataset metadata

Laws of Thermodynamics; Physical Model Grounding Reasoning, Numerical Reasoning

\paragraph{Recorded answer/context.}
D: D: -0.0147 kJ

\paragraph{Figure/image path.}
/root/proposal\_for\_physic/science-mango/phyx\_mini\_run/phyx\_data/test\_image/460.png

\paragraph{Formalization target.}
create a compiling Lean file with sorry bodies at `PhyXMiniProblems/problem\_phyx\_mini\_0460.lean`. The Lean declarations must preserve the physical quantities, dimensions or dimensional roles, figure labels, governing-law hypotheses, and final relation expressed by this problem.
Use Mathlib/Physlib names found through LeanExplore where available. If a physics API is missing, introduce faithful local abstractions rather than scalar placeholder aliases.

\begin{theorem}[Physics formalization target]
\label{thm:physics:phyx_mini_0460:target}
\lean{PhyXMiniProblems.ProblemPhyXMini0460.problem_phyx_mini_0460}
\uses{def:physics:phyx-mini-0460:phyxminiproblems-problemphyxmini0460-energyinjoules, def:physics:phyx-mini-0460:phyxminiproblems-problemphyxmini0460-energyinkilojoules, def:physics:phyx-mini-0460:phyxminiproblems-problemphyxmini0460-airpolytropiccompression, def:physics:phyx-mini-0460:phyxminiproblems-problemphyxmini0460-matchesstatedscenario, def:physics:phyx-mini-0460:phyxminiproblems-problemphyxmini0460-matchesproblemreadouts, def:physics:phyx-mini-0460:phyxminiproblems-problemphyxmini0460-usescoldairstandardcaloricmodel, def:physics:phyx-mini-0460:phyxminiproblems-problemphyxmini0460-matchessuppliedpistoncylinderfigure, def:physics:phyx-mini-0460:phyxminiproblems-problemphyxmini0460-hasphysicalpolytropicparameters, def:physics:phyx-mini-0460:phyxminiproblems-problemphyxmini0460-satisfiespolytropicidealairlaws, def:physics:phyx-mini-0460:phyxminiproblems-problemphyxmini0460-exactmodeledheattransferinjoules, def:physics:phyx-mini-0460:phyxminiproblems-problemphyxmini0460-isuniqueclosestdisplayedheat}
The assigned autoformalize agent should translate this physics problem into Lean declarations in the covered file, with theorem and lemma proof bodies written as `by sorry`.
\end{theorem}
\begin{proof}
This is an autoformalization task, not a proof task. Produce statements that can later be proved by the physics prover without weakening the physical model.
\end{proof}

% --- Archon declaration topology begin ---
\section{Declaration topology}
\begin{definition}[Gas Volume]
  \label{def:physics:phyx-mini-0460:phyxminiproblems-problemphyxmini0460-gasvolume}
  \lean{PhyXMiniProblems.ProblemPhyXMini0460.GasVolume}
  A nonnegative physical gas volume, carrying dimension L³
\end{definition}
\begin{proof}
  This is the declared model definition of Gas Volume.
\end{proof}
\begin{definition}[volume In Cubic Meters]
  \label{def:physics:phyx-mini-0460:phyxminiproblems-problemphyxmini0460-volumeincubicmeters}
  \lean{PhyXMiniProblems.ProblemPhyXMini0460.volumeInCubicMeters}
  \uses{def:physics:phyx-mini-0460:phyxminiproblems-problemphyxmini0460-gasvolume}
  Read a physical gas volume in coherent SI cubic metres
\end{definition}
\begin{proof}
  This is the declared model definition of volume In Cubic Meters.
\end{proof}
\begin{definition}[volume In Liters]
  \label{def:physics:phyx-mini-0460:phyxminiproblems-problemphyxmini0460-volumeinliters}
  \lean{PhyXMiniProblems.ProblemPhyXMini0460.volumeInLiters}
  \uses{def:physics:phyx-mini-0460:phyxminiproblems-problemphyxmini0460-gasvolume, def:physics:phyx-mini-0460:phyxminiproblems-problemphyxmini0460-volumeincubicmeters}
  Read a physical gas volume in litres
\end{definition}
\begin{proof}
  This is the declared model definition of volume In Liters.
\end{proof}
\begin{definition}[pressure In Pascals]
  \label{def:physics:phyx-mini-0460:phyxminiproblems-problemphyxmini0460-pressureinpascals}
  \lean{PhyXMiniProblems.ProblemPhyXMini0460.pressureInPascals}
  Read a physical pressure in pascals
\end{definition}
\begin{proof}
  This is the declared model definition of pressure In Pascals.
\end{proof}
\begin{definition}[pressure In Kilopascals]
  \label{def:physics:phyx-mini-0460:phyxminiproblems-problemphyxmini0460-pressureinkilopascals}
  \lean{PhyXMiniProblems.ProblemPhyXMini0460.pressureInKilopascals}
  \uses{def:physics:phyx-mini-0460:phyxminiproblems-problemphyxmini0460-pressureinpascals}
  Read a physical pressure in kilopascals
\end{definition}
\begin{proof}
  This is the declared model definition of pressure In Kilopascals.
\end{proof}
\begin{definition}[temperature In Kelvins]
  \label{def:physics:phyx-mini-0460:phyxminiproblems-problemphyxmini0460-temperatureinkelvins}
  \lean{PhyXMiniProblems.ProblemPhyXMini0460.temperatureInKelvins}
  Read a Physlib absolute temperature in kelvin
\end{definition}
\begin{proof}
  This is the declared model definition of temperature In Kelvins.
\end{proof}
\begin{definition}[temperature In Degrees Celsius]
  \label{def:physics:phyx-mini-0460:phyxminiproblems-problemphyxmini0460-temperatureindegreescelsius}
  \lean{PhyXMiniProblems.ProblemPhyXMini0460.temperatureInDegreesCelsius}
  \uses{def:physics:phyx-mini-0460:phyxminiproblems-problemphyxmini0460-temperatureinkelvins}
  Read the same temperature on the degrees-Celsius scale
\end{definition}
\begin{proof}
  This is the declared model definition of temperature In Degrees Celsius.
\end{proof}
\begin{definition}[energy In Joules]
  \label{def:physics:phyx-mini-0460:phyxminiproblems-problemphyxmini0460-energyinjoules}
  \lean{PhyXMiniProblems.ProblemPhyXMini0460.energyInJoules}
  Read signed heat, work, or internal energy in joules
\end{definition}
\begin{proof}
  This is the declared model definition of energy In Joules.
\end{proof}
\begin{definition}[energy In Kilojoules]
  \label{def:physics:phyx-mini-0460:phyxminiproblems-problemphyxmini0460-energyinkilojoules}
  \lean{PhyXMiniProblems.ProblemPhyXMini0460.energyInKilojoules}
  \uses{def:physics:phyx-mini-0460:phyxminiproblems-problemphyxmini0460-energyinjoules}
  Read signed heat, work, or internal energy in kilojoules
\end{definition}
\begin{proof}
  This is the declared model definition of energy In Kilojoules.
\end{proof}
\begin{definition}[Endpoint]
  \label{def:physics:phyx-mini-0460:phyxminiproblems-problemphyxmini0460-endpoint}
  \lean{PhyXMiniProblems.ProblemPhyXMini0460.Endpoint}
  The initial and final equilibrium states of the air
\end{definition}
\begin{proof}
  This is the declared model definition of Endpoint.
\end{proof}
\begin{definition}[Air State]
  \label{def:physics:phyx-mini-0460:phyxminiproblems-problemphyxmini0460-airstate}
  \lean{PhyXMiniProblems.ProblemPhyXMini0460.AirState}
  \uses{def:physics:phyx-mini-0460:phyxminiproblems-problemphyxmini0460-gasvolume}
  A macroscopic equilibrium state of the fixed air sample
\end{definition}
\begin{proof}
  This is the declared model definition of Air State.
\end{proof}
\begin{definition}[Working Fluid]
  \label{def:physics:phyx-mini-0460:phyxminiproblems-problemphyxmini0460-workingfluid}
  \lean{PhyXMiniProblems.ProblemPhyXMini0460.WorkingFluid}
  Working-fluid identity named by the problem statement and figure
\end{definition}
\begin{proof}
  This is the declared model definition of Working Fluid.
\end{proof}
\begin{definition}[Gas Model]
  \label{def:physics:phyx-mini-0460:phyxminiproblems-problemphyxmini0460-gasmodel}
  \lean{PhyXMiniProblems.ProblemPhyXMini0460.GasModel}
  Equation-of-state model used for the air
\end{definition}
\begin{proof}
  This is the declared model definition of Gas Model.
\end{proof}
\begin{definition}[Inventory Regime]
  \label{def:physics:phyx-mini-0460:phyxminiproblems-problemphyxmini0460-inventoryregime}
  \lean{PhyXMiniProblems.ProblemPhyXMini0460.InventoryRegime}
  Whether mass crosses the piston--cylinder boundary
\end{definition}
\begin{proof}
  This is the declared model definition of Inventory Regime.
\end{proof}
\begin{definition}[Mechanical Regime]
  \label{def:physics:phyx-mini-0460:phyxminiproblems-problemphyxmini0460-mechanicalregime}
  \lean{PhyXMiniProblems.ProblemPhyXMini0460.MechanicalRegime}
  Mechanical regime of the moving piston
\end{definition}
\begin{proof}
  This is the declared model definition of Mechanical Regime.
\end{proof}
\begin{definition}[Process Direction]
  \label{def:physics:phyx-mini-0460:phyxminiproblems-problemphyxmini0460-processdirection}
  \lean{PhyXMiniProblems.ProblemPhyXMini0460.ProcessDirection}
  Direction of the volume change
\end{definition}
\begin{proof}
  This is the declared model definition of Process Direction.
\end{proof}
\begin{definition}[Figure Region]
  \label{def:physics:phyx-mini-0460:phyxminiproblems-problemphyxmini0460-figureregion}
  \lean{PhyXMiniProblems.ProblemPhyXMini0460.FigureRegion}
  Regions visually distinguished in the supplied piston--cylinder image
\end{definition}
\begin{proof}
  This is the declared model definition of Figure Region.
\end{proof}
\begin{definition}[Piston Outline]
  \label{def:physics:phyx-mini-0460:phyxminiproblems-problemphyxmini0460-pistonoutline}
  \lean{PhyXMiniProblems.ProblemPhyXMini0460.PistonOutline}
  Line style used to distinguish the two piston positions in the image
\end{definition}
\begin{proof}
  This is the declared model definition of Piston Outline.
\end{proof}
\begin{definition}[Horizontal Placement]
  \label{def:physics:phyx-mini-0460:phyxminiproblems-problemphyxmini0460-horizontalplacement}
  \lean{PhyXMiniProblems.ProblemPhyXMini0460.HorizontalPlacement}
  Horizontal ordering of the two depicted piston positions
\end{definition}
\begin{proof}
  This is the declared model definition of Horizontal Placement.
\end{proof}
\begin{definition}[Piston Cylinder Figure]
  \label{def:physics:phyx-mini-0460:phyxminiproblems-problemphyxmini0460-pistoncylinderfigure}
  \lean{PhyXMiniProblems.ProblemPhyXMini0460.PistonCylinderFigure}
  \uses{def:physics:phyx-mini-0460:phyxminiproblems-problemphyxmini0460-endpoint, def:physics:phyx-mini-0460:phyxminiproblems-problemphyxmini0460-figureregion, def:physics:phyx-mini-0460:phyxminiproblems-problemphyxmini0460-pistonoutline, def:physics:phyx-mini-0460:phyxminiproblems-problemphyxmini0460-horizontalplacement}
  Qualitative data carried by the supplied piston--cylinder schematic
\end{definition}
\begin{proof}
  This is the declared model definition of Piston Cylinder Figure.
\end{proof}
\begin{definition}[Air Polytropic Compression]
  \label{def:physics:phyx-mini-0460:phyxminiproblems-problemphyxmini0460-airpolytropiccompression}
  \lean{PhyXMiniProblems.ProblemPhyXMini0460.AirPolytropicCompression}
  \uses{def:physics:phyx-mini-0460:phyxminiproblems-problemphyxmini0460-endpoint, def:physics:phyx-mini-0460:phyxminiproblems-problemphyxmini0460-airstate, def:physics:phyx-mini-0460:phyxminiproblems-problemphyxmini0460-workingfluid, def:physics:phyx-mini-0460:phyxminiproblems-problemphyxmini0460-gasmodel, def:physics:phyx-mini-0460:phyxminiproblems-problemphyxmini0460-inventoryregime, def:physics:phyx-mini-0460:phyxminiproblems-problemphyxmini0460-mechanicalregime, def:physics:phyx-mini-0460:phyxminiproblems-problemphyxmini0460-processdirection, def:physics:phyx-mini-0460:phyxminiproblems-problemphyxmini0460-pistoncylinderfigure}
  The complete physical process, including unknown heat and work
\end{definition}
\begin{proof}
  This is the declared model definition of Air Polytropic Compression.
\end{proof}
\begin{definition}[Matches Stated Scenario]
  \label{def:physics:phyx-mini-0460:phyxminiproblems-problemphyxmini0460-matchesstatedscenario}
  \lean{PhyXMiniProblems.ProblemPhyXMini0460.MatchesStatedScenario}
  \uses{def:physics:phyx-mini-0460:phyxminiproblems-problemphyxmini0460-airpolytropiccompression}
  Qualitative physical information stated in the problem
\end{definition}
\begin{proof}
  This is the declared model definition of Matches Stated Scenario.
\end{proof}
\begin{definition}[Matches Problem Readouts]
  \label{def:physics:phyx-mini-0460:phyxminiproblems-problemphyxmini0460-matchesproblemreadouts}
  \lean{PhyXMiniProblems.ProblemPhyXMini0460.MatchesProblemReadouts}
  \uses{def:physics:phyx-mini-0460:phyxminiproblems-problemphyxmini0460-volumeincubicmeters, def:physics:phyx-mini-0460:phyxminiproblems-problemphyxmini0460-volumeinliters, def:physics:phyx-mini-0460:phyxminiproblems-problemphyxmini0460-pressureinkilopascals, def:physics:phyx-mini-0460:phyxminiproblems-problemphyxmini0460-temperatureindegreescelsius, def:physics:phyx-mini-0460:phyxminiproblems-problemphyxmini0460-airpolytropiccompression}
  Exact transcription of the scalar data printed in the question
\end{definition}
\begin{proof}
  This is the declared model definition of Matches Problem Readouts.
\end{proof}
\begin{definition}[Uses Cold Air Standard Caloric Model]
  \label{def:physics:phyx-mini-0460:phyxminiproblems-problemphyxmini0460-usescoldairstandardcaloricmodel}
  \lean{PhyXMiniProblems.ProblemPhyXMini0460.UsesColdAirStandardCaloricModel}
  \uses{def:physics:phyx-mini-0460:phyxminiproblems-problemphyxmini0460-airpolytropiccompression}
  The cold-air-standard caloric index is model input, separated from the values explicitly printed in the question. It is not a heat-transfer conclusion
\end{definition}
\begin{proof}
  This is the declared model definition of Uses Cold Air Standard Caloric Model.
\end{proof}
\begin{definition}[Matches Supplied Piston Cylinder Figure]
  \label{def:physics:phyx-mini-0460:phyxminiproblems-problemphyxmini0460-matchessuppliedpistoncylinderfigure}
  \lean{PhyXMiniProblems.ProblemPhyXMini0460.MatchesSuppliedPistonCylinderFigure}
  \uses{def:physics:phyx-mini-0460:phyxminiproblems-problemphyxmini0460-airpolytropiccompression}
  Figure-derived label, casing, and relative piston-position information
\end{definition}
\begin{proof}
  This is the declared model definition of Matches Supplied Piston Cylinder Figure.
\end{proof}
\begin{definition}[Has Physical Polytropic Parameters]
  \label{def:physics:phyx-mini-0460:phyxminiproblems-problemphyxmini0460-hasphysicalpolytropicparameters}
  \lean{PhyXMiniProblems.ProblemPhyXMini0460.HasPhysicalPolytropicParameters}
  \uses{def:physics:phyx-mini-0460:phyxminiproblems-problemphyxmini0460-volumeincubicmeters, def:physics:phyx-mini-0460:phyxminiproblems-problemphyxmini0460-pressureinpascals, def:physics:phyx-mini-0460:phyxminiproblems-problemphyxmini0460-temperatureinkelvins, def:physics:phyx-mini-0460:phyxminiproblems-problemphyxmini0460-airpolytropiccompression}
  Positivity and nondegeneracy conditions for the physical process
\end{definition}
\begin{proof}
  This is the declared model definition of Has Physical Polytropic Parameters.
\end{proof}
\begin{definition}[Satisfies Polytropic Ideal Air Laws]
  \label{def:physics:phyx-mini-0460:phyxminiproblems-problemphyxmini0460-satisfiespolytropicidealairlaws}
  \lean{PhyXMiniProblems.ProblemPhyXMini0460.SatisfiesPolytropicIdealAirLaws}
  \uses{def:physics:phyx-mini-0460:phyxminiproblems-problemphyxmini0460-volumeincubicmeters, def:physics:phyx-mini-0460:phyxminiproblems-problemphyxmini0460-pressureinpascals, def:physics:phyx-mini-0460:phyxminiproblems-problemphyxmini0460-temperatureinkelvins, def:physics:phyx-mini-0460:phyxminiproblems-problemphyxmini0460-energyinjoules, def:physics:phyx-mini-0460:phyxminiproblems-problemphyxmini0460-airpolytropiccompression}
  These are general endpoint laws for a calorically perfect ideal gas on a quasi-equilibrium polytropic path. None contains the requested heat value or an answer-choice label
\end{definition}
\begin{proof}
  This is the declared model definition of Satisfies Polytropic Ideal Air Laws.
\end{proof}
\begin{definition}[exact Modeled Heat Transfer In Joules]
  \label{def:physics:phyx-mini-0460:phyxminiproblems-problemphyxmini0460-exactmodeledheattransferinjoules}
  \lean{PhyXMiniProblems.ProblemPhyXMini0460.exactModeledHeatTransferInJoules}
  The unrounded cold-air-standard result in joules. It is obtained by adding ΔU = Δ(PV)/(γ - 1) and W = Δ(PV)/(1 - n), with Δ(PV) = P₁ V₁ (6\textasciicircum{}(n-1) - 1)
\end{definition}
\begin{proof}
  This is the declared model definition of exact Modeled Heat Transfer In Joules.
\end{proof}
\begin{definition}[Answer Choice]
  \label{def:physics:phyx-mini-0460:phyxminiproblems-problemphyxmini0460-answerchoice}
  \lean{PhyXMiniProblems.ProblemPhyXMini0460.AnswerChoice}
  Labels printed beside the four candidate heat-transfer values
\end{definition}
\begin{proof}
  This is the declared model definition of Answer Choice.
\end{proof}
\begin{definition}[displayed Heat Transfer In Kilojoules]
  \label{def:physics:phyx-mini-0460:phyxminiproblems-problemphyxmini0460-displayedheattransferinkilojoules}
  \lean{PhyXMiniProblems.ProblemPhyXMini0460.displayedHeatTransferInKilojoules}
  \uses{def:physics:phyx-mini-0460:phyxminiproblems-problemphyxmini0460-answerchoice}
  Displayed candidate values, interpreted in kilojoules with heat into air positive
\end{definition}
\begin{proof}
  This is the declared model definition of displayed Heat Transfer In Kilojoules.
\end{proof}
\begin{definition}[recorded Dataset Answer]
  \label{def:physics:phyx-mini-0460:phyxminiproblems-problemphyxmini0460-recordeddatasetanswer}
  \lean{PhyXMiniProblems.ProblemPhyXMini0460.recordedDatasetAnswer}
  \uses{def:physics:phyx-mini-0460:phyxminiproblems-problemphyxmini0460-answerchoice}
  The answer label recorded by the source dataset, retained as metadata only
\end{definition}
\begin{proof}
  This is the declared model definition of recorded Dataset Answer.
\end{proof}
\begin{definition}[Is Unique Closest Displayed Heat]
  \label{def:physics:phyx-mini-0460:phyxminiproblems-problemphyxmini0460-isuniqueclosestdisplayedheat}
  \lean{PhyXMiniProblems.ProblemPhyXMini0460.IsUniqueClosestDisplayedHeat}
  \uses{def:physics:phyx-mini-0460:phyxminiproblems-problemphyxmini0460-answerchoice, def:physics:phyx-mini-0460:phyxminiproblems-problemphyxmini0460-displayedheattransferinkilojoules}
  A displayed option is strictly closer to a modeled heat than every alternative
\end{definition}
\begin{proof}
  This is the declared model definition of Is Unique Closest Displayed Heat.
\end{proof}
\begin{lemma}[heat Transfer eq exact Modeled Value]
  \label{lem:physics:phyx-mini-0460:phyxminiproblems-problemphyxmini0460-heattransfer-eq-exactmodeledvalue}
  \lean{PhyXMiniProblems.ProblemPhyXMini0460.heatTransfer_eq_exactModeledValue}
  \uses{def:physics:phyx-mini-0460:phyxminiproblems-problemphyxmini0460-energyinjoules, def:physics:phyx-mini-0460:phyxminiproblems-problemphyxmini0460-airpolytropiccompression, def:physics:phyx-mini-0460:phyxminiproblems-problemphyxmini0460-matchesstatedscenario, def:physics:phyx-mini-0460:phyxminiproblems-problemphyxmini0460-matchesproblemreadouts, def:physics:phyx-mini-0460:phyxminiproblems-problemphyxmini0460-usescoldairstandardcaloricmodel, def:physics:phyx-mini-0460:phyxminiproblems-problemphyxmini0460-hasphysicalpolytropicparameters, def:physics:phyx-mini-0460:phyxminiproblems-problemphyxmini0460-satisfiespolytropicidealairlaws, def:physics:phyx-mini-0460:phyxminiproblems-problemphyxmini0460-exactmodeledheattransferinjoules}
  The physical laws determine the unrounded heat-transfer value. This is a derived conclusion, not a premise of the model
\end{lemma}
\begin{proof}
  The conclusion follows from the governing relations and figure readouts named in the statement of heat Transfer eq exact Modeled Value.
\end{proof}
\begin{lemma}[exact Modeled Heat Transfer selects choice D]
  \label{lem:physics:phyx-mini-0460:phyxminiproblems-problemphyxmini0460-exactmodeledheattransfer-selects-choice-d}
  \lean{PhyXMiniProblems.ProblemPhyXMini0460.exactModeledHeatTransfer_selects_choice_D}
  \uses{def:physics:phyx-mini-0460:phyxminiproblems-problemphyxmini0460-exactmodeledheattransferinjoules, def:physics:phyx-mini-0460:phyxminiproblems-problemphyxmini0460-isuniqueclosestdisplayedheat}
  The unrounded cold-air model selects the printed value -0.0147 kJ
\end{lemma}
\begin{proof}
  The conclusion follows from the governing relations and figure readouts named in the statement of exact Modeled Heat Transfer selects choice D.
\end{proof}
% --- Archon declaration topology end ---
% --- Archon physics formalization source end ---
